\newpage
\chapter{KESIMPULAN DAN SARAN} \label{Bab V}

\section{Kesimpulan} \label{V.Kesimpulan}
Berdasarkan hasil penelitian, pemodelan, dan evaluasi yang telah dilakukan terhadap algoritma \textit{Random Forest} dan interpretasi SHAP untuk menganalisis determinan \textit{stunting} tingkat provinsi di Indonesia menggunakan data SSGI 2024, dapat ditarik kesimpulan sebagai berikut:

\begin{enumerate}
    \item Penelitian ini berhasil mengidentifikasi faktor-faktor determinan utama penyebab prevalensi \textit{stunting} tingkat provinsi di Indonesia. Berdasarkan pemeringkatan fitur pada model, variabel keberagaman konsumsi makanan balita (seperti ketiadaan MDD dan MMFF), kondisi kehamilan (usia ibu berisiko), akses infrastruktur dasar (sanitasi berisiko), serta cakupan bantuan gizi (ketiadaan PMT) menjadi prediktor paling dominan. Temuan ini menegaskan bahwa masalah malnutrisi komunal di Indonesia sangat dipengaruhi oleh kombinasi antara ketahanan pangan harian, kesiapan fisik ibu, dan kebersihan lingkungan tempat tinggal.
    
    \item Algoritma \textit{Random Forest Regression} terbukti tangguh dan mumpuni dalam memprediksi tingkat prevalensi \textit{stunting} antarprovinsi. Model ini mampu memetakan pola data agregat makro yang bersifat non-linear dengan tingkat akurasi yang sangat baik, dibuktikan dengan perolehan metrik evaluasi regresi berupa nilai $R^2$ sebesar 0,9090, \textit{Mean Absolute Error} (MAE) sebesar 1,40, \textit{Root Mean Squared Error} (RMSE) sebesar 1,77, serta tingkat kesalahan prediksi (\textit{Mean Absolute Percentage Error}/MAPE) sebesar 7,22\%. Kinerja ini menunjukkan bahwa algoritma \textit{Random Forest} sangat layak dan kompetitif untuk digunakan pada pemodelan data survei kesehatan masyarakat berskala nasional.
  
    \item Penerapan metode interpretabilitas SHAP (\textit{SHapley Additive exPlanations}) berhasil membuka \textit{black-box} pada model \textit{Random Forest} dengan mengukur kontribusi spesifik setiap fitur. Melalui pemetaan \textit{Dependence Plot}, model mengekstraksi ambang batas toleransi kritis komunal secara matematis, di antaranya: batas ketiadaan keberagaman makanan pokok (MDD) di angka 55\%, batas ketiadaan asupan susu (MMFF) di angka 22\%, proporsi kehamilan berisiko di batas 14\%, serta ketiadaan sanitasi layak di angka 10\%. Evaluasi arah fitur membuktikan bahwa suatu provinsi akan mengalami lonjakan risiko prevalensi secara drastis apabila persentase ketiadaan fasilitas atau gizi melampaui titik ambang batas tersebut.
\end{enumerate}

\section{Saran} \label{V.Saran}
Berdasarkan hasil analisis serta keterbatasan yang terdapat dalam penelitian ini, terdapat beberapa saran yang dapat dijadikan sebagai bahan pengembangan pada penelitian berikutnya maupun dalam penerapan secara praktis, yaitu sebagai berikut:

\begin{enumerate}
    \item Penelitian ini masih menggunakan data agregat pada tingkat provinsi sehingga hasil analisis yang diperoleh lebih menggambarkan kondisi secara umum. Oleh karena itu, penelitian selanjutnya disarankan untuk menggunakan data dengan cakupan wilayah yang lebih spesifik, seperti tingkat kabupaten/kota. Penggunaan data yang lebih rinci diharapkan dapat memberikan gambaran faktor penyebab \textit{stunting} secara lebih mendalam sesuai karakteristik masing-masing daerah.
    
    \item Penelitian ini menggunakan metode SHAP untuk melakukan interpretasi terhadap model \textit{Random Forest} yang bersifat \textit{black-box}. Pada penelitian selanjutnya, disarankan untuk melakukan perbandingan dengan pendekatan \textit{Explainable AI} (XAI) lainnya, seperti LIME (\textit{Local Interpretable Model-Agnostic Explanations}). Perbandingan antar metode interpretasi tersebut dapat digunakan untuk melihat konsistensi hasil analisis, khususnya pada penentuan titik ambang batas dari determinan \textit{stunting}, sehingga hasil yang diperoleh memiliki dasar analisis komputasional yang lebih kuat dan meyakinkan.
    
    \item Hasil penelitian ini juga menunjukkan adanya beberapa titik ambang batas penting, seperti persentase usia ibu berisiko sebesar 14\% dan sanitasi buruk sebesar 10\%. Nilai tersebut dapat dimanfaatkan sebagai acuan awal bagi pemangku kebijakan dalam mengevaluasi efektivitas program penanganan \textit{stunting}, khususnya yang berkaitan dengan perbaikan gizi dan sanitasi lingkungan.
\end{enumerate}